%% kapa: exam=1 round=27 year=2016 period=1 form=A subject=civil_law subject_ko=민법 seq=21 number=21
%% source: Q-Net 감정평가사 기출문제 (cst00309) artlSeq=5169402 file=extracted/1008/5169402/제27회 감정평가사 1차 문제/제27회 감정평가사 1교시문제(A형).pdf
%% conversion: kordoc 4.24.51 -> tools/gen1.py (auto); review=proofread-sonnet; flags=none
\begin{question}{21}
가등기에 관한 설명으로 옳지 않은 것은? (다툼이 있으면 판례에 따름)
\begin{choices}
  \item 가등기상의 권리의 이전등기를 가등기에 대한 부기등기의 형식으로 할 수 있다.
  \item 저당권설정등기청구권을 보전하기 위한 가등기는 인정되지 않는다.
  \item 가등기에 기한 본등기가 경료되더라도 본등기에 의한 물권변동의 효력이 가등기한 때로 소급하여 발생하는 것은 아니다.
  \item 가등기가 부적법하게 말소된 후 소유권이전등기를 마친 제3자는 가등기의 회복등기절차에서 승낙의무가 있다.
  \item 가등기에 기한 본등기청구권과 소유권이전등기청구권은 그 등기원인이 동일하더라도 서로 다른 청구권으로 보아야 한다.
\end{choices}
\end{question}
\answer{2}
